\section{Выводы}

В результате выполнения лабораторной работы был успешно реализован алгоритм точного поиска образца в тексте на основе Z-функции.

\textbf{Об алгоритме.} На практике было доказано, что классические алгоритмы на строках легко обобщаются. Z-функция, изначально спроектированная для работы с массивами символов (строками), без изменения базовой логики $O(N)$ была перенесена на работу с массивами слов. Понятие \enquote{символ} было заменено на структуру данных \texttt{std::string}, а равенство символов — на перегруженный оператор \texttt{==}.

\textbf{О вводе-выводе.} Одной из главных технических сложностей оказалось сохранение метаданных (номеров строк и слов) при чтении из потока. Стандартный оператор \texttt{std::cin >>} игнорирует пробельные символы и переносы каретки (\texttt{\textbackslash n}), из-за чего невозможно было отследить переход на новую строку. Использование посимвольного чтения через \texttt{std::cin.get()} с реализацией собственного простейшего лексического анализатора решило эту проблему и позволило in-place (без выделения дополнительной памяти) осуществлять приведение букв к нижнему регистру.

\textbf{О производительности.} В вырожденных ситуациях (худший случай), где наивный подход и стандартные библиотечные функции \enquote{зависают} из-за квадратичной деградации $O(N \times M)$, Z-функция обеспечивает стабильную линейную скорость. Это доказало критическую важность выбора правильного алгоритма при проектировании высоконагруженных систем обработки текстов.
\pagebreak